% Nouveau EGA V (cote 25): archive cote 25, site page 45 (PDF page 46).
% typescript. Draft compared: none.
% AI-transcribed by Claude Opus 5.5 (claude-opus-5-5, adaptive thinking, effort high), Message Batches API, batch ega25b (msgbatch_01A7n5YuXubUPNK1zsepi8RG), 2026-11-13; original AI transcription, subsequently checked against both supplied scan parts; see PAGES_CHECKED.tsv and CHECK_DECISIONS.tsv.
\scanpage{25/45}
% NOTE: The small left-margin blot has uncertain function. The detached acute-like pen stroke over lesquels is retained without assigning grammatical force. The circled insertion places sing on a lower baseline than sup.

\mspage{\hand{41}}

hg$_1$ :  $g^{-1}(U) \hand{\longrightarrow}$ PxT  n'est autre que l'inclusion canonique, puisqu'il
en est ainsi de sa restriction à $g^{-1}(V) \hand{\rightarrow}$ V.  
 Il en résulte que h  se factorise
à travers l'adhérence schématique  $\overset{\hand{\rule{1.1em}{0.35pt}}}{Y^1}$ de Y$^1$ dans  PxT , donc que
l'image inverse de Y$^1$ (qui est un ouvert de ladite adhérence) par h
est un ouvert de U, soit U' . En vertu de  h g$_1$ = inclusion, on
voit alors aussitôt que  g$_1$ induit un isomorphisme  $g^{-1}(U) \hand{\simeq}$ U', et
 cela montre bien que  $g^{-1}(U) \hand{\longrightarrow}$ U est une immersion ouverte.
Lorsque X est propre sur k i.e. f ermé dans P, alors  g est propre donc
surjectif, donc  g$_1$ : $g^{-1}(U) \hand{\longrightarrow}$ U est surjectif, donc un isomorphisme.
D'ailleurs, si  W est un ouvert de T tel que  $g^{-1}(W) \hand{\longrightarrow}$ W soit un isomor%
phisme, il s'ensuit que W est lisse puisque Y$^1$ est lisse, donc W$\hand{\subset}$U.
Cela achève la démonstration de 8.9. 

  Les assertion\hand{s} finale\hand{s} de 8.8.\hand{,} $\underline{Y}^{\text{sup}}\text{sing}=\emptyset$ ou dim $\underline{Y}^{\text{sup sing}}$ \unsure{$\leq$} r-2\hand{,} et
dim T = r-1 , sont triviales;  la première résulte d\unsure{e} 
  fait que $\underline{Y}^{\text{sing}}$, de dim r  est irréductible, et du fait que
 $\underline{Y}^{\text{sup sing}}$  est défini par l'annulation d'une section D d'un Module inversible ; la deuxième du fait que, L étant fini sur K, on a
deg tr L/k = deg tr K/k, i.e.   dim T = dim $\underline{Y}^{\text{sing}}$ = r-1.

\uline{Remarque} 8.10.  Comme on a signalé dans  8.9. , avec les notations de
ce corollaire, on a  $g^{-1}(U) \hand{\subset} \underline{Y}^{\text{sing}} - \underline{Y}^{\text{sup sing}}$ ; mais on notera que
même si X est fermé dans P,  cette inclusion n'est pas nécessairement
\placedmargin{\hand{\unsure{$\cdot$}}} une égalité, en d'autres termes (notant que $g^{-1}(U)$ n'est autre que
l'ensemble des points où g est \uline{étale}, alors que $\underline{Y}^{\text{sing}}$-$\underline{Y}^{\text{sup sing}}$ est
lèensemble des points où g est \uline{non ramifié}) il peut y avoir des points
y de $\underline{Y}^{\text{sing}}$ en \hand{´}lesquels  g est  non ramifié, mais non étale (ce qui
 implique d'ailleurs que g(y) est un point non  géom normal et même
non géom unibranche de  T)\hand{.}  En termes géométriques, cela corres%
pond au phénomène suivant: on peut avoir un hyperplan tangent non oscul
% NOTE: the last line is cut at the right edge of the scan
\placedmargin{\hand{45}}
